\section{What the search is good for}\label{sec:predict}

The preceding sections are negative about \emph{finality}. Here we explain why the
search for compact, unified theories is still worthwhile. The reason is prediction. We
first state the strongest predictive guarantee available, then show that no usable
method has it, and finally say what compression offers in its place.

\paragraph{The ideal benchmark.}
Take stochastic models as in \cref{def:model}(b), with the enumeration
$\nu_1,\nu_2,\ldots$ fixed there and the outcomes of a record forming a string
$z=z_1\cdots z_n$, and write $\K(\nu)=\min\{\K(i):\nu_i=\nu\}$ for
the least prefix complexity of an index of $\nu$. The universal mixture
\[
  \Univ(z)=\sum_{i\ge1}2^{-\K(i)}\nu_i(z)
\]
weights every model by its compressed length and predicts with
$\Univ(a\mid z)=\Univ(za)/\Univ(z)$. Up to a multiplicative constant it is Solomonoff's
a priori probability \citep{solomonoff1978,li-vitanyi2019}.

\begin{theorem}[Relative optimality]\label{thm:pred-relative}
For every lower semicomputable semimeasure $\nu$ and every string $z=z_1\cdots z_n$,
\[
  \sum_{t=1}^{n}-\log_2\Univ(z_t\mid z_{<t})\;\le\;-\log_2\nu(z)+\K(\nu).
\]
\end{theorem}

\begin{proof}
If $\nu(z)=0$ the bound is trivial. Otherwise, for every index $i$ of $\nu$ the mixture
contains the term $2^{-\K(i)}\nu(z)$, so $\Univ(z)\ge2^{-\K(\nu)}\nu(z)$
\citep[Ch.~4]{li-vitanyi2019}. The conditional
log-losses telescope to $-\log_2\Univ(z)+\log_2\Univ(\varepsilon)\le-\log_2\Univ(z)$.
\end{proof}

The bound holds for every sequence, however it was produced, and it assumes nothing
about the source. Its content is comparative: on any data, the mixture loses at most
$\K(\nu)$ bits more than any given model $\nu$. The constant depends on the model and
on the reference machine.

\begin{proposition}[No computable predictor dominates]\label{pred:prop:no-computable}
Let $\rho$ be a predictor with computable conditionals satisfying
$\sum_a\rho(a\mid z)\le1$. Write $\K(\rho)$ for the least prefix complexity of a program computing $\rho$'s
conditionals. There is a computable sequence $\omega$ on which $\rho$'s
cumulative log-loss after $n$ symbols is at least $n\log_2|\Sig|-O(1)$. Meanwhile the
deterministic model of $\omega$ has log-loss $0$ and complexity at most
$\K(\rho)+O(1)$. Hence no computable predictor satisfies a bound of the form of
\cref{thm:pred-relative}.
\end{proposition}

\begin{proof}
At each step, approximate $\rho(a\mid\omega_{<t})$ for every $a\in\Sig$ to within
$2^{-t}/|\Sig|$ and let $\omega_t$ be the symbol with the least approximation, breaking
ties in a fixed order. The least conditional is at most $1/|\Sig|$, so
$\rho(\omega_t\mid\omega_{<t})\le(1+2^{-t+1})/|\Sig|$. Summing $-\log_2$ of these gives
the bound, since $\sum_t\log_2(1+2^{-t+1})<\infty$. The construction is uniform in a
program for $\rho$.
\end{proof}

By \cref{prem:effective}, then, the benchmark is not a usable method: every usable
predictor is computable, and for each there is a computable sequence on which it does as
badly as possible. What \cref{thm:pred-relative} supplies is a sense in which weighting
models by their compressed length is an optimal bias up to $K(\nu)$: relative to every lower
semicomputable model, on every sequence, it loses at most a constant. It is an ideal;
usable approximations aim at it, and each is defeated on some computable sequence. Martins draws a related moral for physics from normative
inductive reasoning: theories should be weighed, not chosen and discarded
\citep{martins2017}.

\paragraph{What compression offers.}
The benchmark nevertheless shows what compression is good for. Every model, a lookup
table included, has outputs at settings not yet observed; what distinguishes a short
model is its weight. The universal mixture gives a model of length $\K(\nu)$ weight at
least $2^{-\K(\nu)}$, so by \cref{thm:pred-relative} predicting as that model does costs
at most $\K(\nu)$ bits, while a lookup table pays about the length of the record.
Compressing what is known therefore commits a theory cheaply to predictions about what
is not yet known, and when those predictions are tested, the bias that chose the theory
is put at risk.
\begin{itemize}
  \item Maxwell's equations unified the known laws of electricity and magnetism. They
        then predicted electromagnetic waves traveling at the speed of light, which
        Hertz produced.
  \item Dirac's equation reconciled quantum mechanics with special relativity for the
        electron. It predicted a particle of the same mass and opposite charge, and the
        positron was found four years later.
  \item Compression can also fail. Le Verrier explained the anomalous precession of
        Mercury, within Newtonian mechanics, by an unseen inner planet, Vulcan. The
        planet does not exist, and the anomaly was explained by general relativity.
\end{itemize}
This is the value of unification, and it is why the search for compact theories of
everything is worth pursuing: a more compact theory of more phenomena makes cheaper
commitments beyond the record, and so it can be tested more severely.

Predictive success does not certify finality. Every success is a finite record, and by
\cref{thm:k1-certify} no finite record certifies that a model is final. Newtonian
mechanics was confirmed by predictions for two centuries, including the discovery of
Neptune, and it was not final. A theory of everything confirmed by every test anyone
performs would be in the same position.
